%!TEX root = ../main.tex
\newpage
\section{Feature Engineering}\label{sec:features_creating}

The goal of the feature engineering phase is to translate each article into a set of quantitative attributes that reflect its journalistic quality. We design these features using principles from computational linguistics and journalism theory to capture various aspects of writing and content. This allows the classification model to assess quality based on interpretable characteristics of the text. (Certain implementation details are omitted here due to confidentiality, but the underlying methods and rationale are described.) The features chosen are grounded in four key dimensions of journalistic quality – \textbf{Diversity}, \textbf{Transparency}, \textbf{Attractiveness}, and \textbf{Independence} – ensuring that our approach aligns with established editorial standards and ethical AI practices. In other words, each feature group corresponds to one or more of these quality dimensions, so that the model’s decisions are based on meaningful, fair criteria rather than arbitrary signals. The following subsections detail each feature group and its relation to the quality dimensions.

\subsection{Linguistic Complexity Features (Attractiveness)}
One important aspect of article quality is the richness and complexity of its language. Linguistic complexity features quantify the sophistication of writing, which corresponds to the Attractiveness dimension of quality – an engaging, well-crafted article typically uses varied and well-structured language. In the Czech language context, we pay special attention to syntax, vocabulary variety, and morphology. This group includes:
\begin{itemize}
    \item \textbf{Syntactic Complexity:} Measures of sentence structure complexity, such as average sentence length and the use of subordinate clauses or compound sentences. High-quality articles often contain longer, information-rich sentences or a mix of sentence lengths that provide detail and nuance. In contrast, lower-quality pieces may rely on very short, simple sentences. We capture this by computing the average number of words per sentence and identifying complex sentence constructions. (For Czech, we note that free word order can make sentences longer or differently structured without losing clarity, so we focus on the presence of well-formed clauses and punctuation rather than word order alone.)
    \item \textbf{Lexical Diversity:} Metrics that assess the variety of vocabulary used, most notably the type-token ratio – the proportion of unique words (types) to total words (tokens) in the article. A higher unique word ratio indicates a broader vocabulary and more varied word choice, which is often a sign of a more informative and engaging article. As shown by the data analysis in Section~\ref{sec:data_analysis}, higher-quality articles tend to employ more unique words than lower-quality ones. For instance, Figure~\ref{fig:num_unique_tokens} in Section~\ref{sec:data_analysis} illustrates that \ac{Class 3} often contain well over 1,000 unique tokens, whereas \ac{Class 0} have far fewer unique words. This confirms that better articles use a wider vocabulary, though there is still some overlap between classes.
    \item \textbf{Morphological Richness:} Features that capture the range of word forms and complex words. Czech is a highly inflected language, meaning words change form (through prefixes, suffixes, case endings, etc.) and new words can be formed via derivation and compounding. We account for derived word forms and compound nouns usage in the text. A quality article might use more varied word forms (indicating precise expression and nuanced meaning) and appropriate compounds, reflecting a command of language. In summary, a high morphological variety signals that the author is leveraging the language’s full expressive capability, aligning with an attractive writing style.
\end{itemize}

Together, these linguistic complexity features describe how sophisticated and varied the article’s language is. They tie into the Attractiveness dimension because an article that is written with complex sentence structures, a broad vocabulary, and rich morphology is generally more engaging and indicative of higher journalistic effort. However, we also remain mindful that clarity is important – complexity should not come at the expense of understanding. (This balance is addressed by the readability features, described later.) Notably, initial data exploration showed that these features are promising: for instance, articles that editors rated as high quality exhibited significantly longer sentences and greater vocabulary variety than low-quality ones, which justifies including these metrics.

\subsection{Structural Composition Features (Transparency)}

Another key aspect of quality is how the article is organized and whether it adheres to good journalistic practices in its structure and use of sources. Structural composition features describe the layout and content elements of an article beyond just the raw text. These features align mostly with the Transparency dimension (and also support Independence), as they reflect how openly and clearly the article presents information and sources. High-quality journalism is usually well-structured and transparent about where information comes from. We include several features in this group:
\begin{itemize}
  \item \textbf{Article Length:} The total length of the article (measured in words). This is a fundamental structural attribute – while not a quality guarantee on its own, length often correlates with depth of content. In Section~\ref{sec:data_analysis}, we observed that better-quality articles tend to be much longer on average. Figure~\ref{fig:num_words_per_label} in Section~\ref{sec:data_analysis} shows the distribution of article lengths by quality class: High-Quality Articles have a wide length distribution with many articles over 1,500 words (and some even above 5,000 words), whereas \ac{Class 0} are mostly under 600 words. This indicates that longer texts are more likely to include detailed explanations, background, and evidence – traits of high-quality reporting.
  \item \textbf{Section Organization:} Features capturing how the article is structured into sections or paragraphs. A well-structured news article typically has a clear introduction (lead), a body divided into coherent paragraphs or sub-sections, and sometimes a conclusion. We quantify this by looking at the number of paragraphs or explicit sections, the presence of subheadings or section titles, and the overall flow. For example, an article with many one-sentence paragraphs or a jumbled structure might indicate lower quality or a rushed piece, whereas one with logically separated sections (e.g. an introduction, several thematic paragraphs, and a wrap-up) suggests careful composition. This structural awareness contributes to transparency and readability – an article that is organized logically is easier for readers to follow and shows that the author/editors put thought into presenting the information clearly.
  \item \textbf{Multimedia Content Metadata:}  Considers the presence and originality of multimedia elements like images, videos, and infographics. Original, journalist-taken photographs significantly enhance transparency, distinguishing higher-quality journalism from stock images. Additionally, metadata such as image descriptions or alt-text indicate editorial commitment to accessibility and clarity.
  \item \textbf{Citations and Source References:} Metrics evaluating how the article cites sources or includes references. A transparent article will often mention where information comes from – for example, quoting officials, referencing reports, or providing hyperlinks to source documents. We also consider source diversity in a basic way: if an article cites several different people or organizations, it’s likely providing a more balanced, independent view (related to the Independence dimension as well). This encourages the classifier to reward articles that follow the journalistic norm of attributing information, aligning with higher quality. Also we classify quotes by origin—whether sourced directly by the journalist (original quotes), from agencies, directly spoken to the channel, or indirectly through platforms like social media (Instagram). Clearly attributed original quotes signal high transparency and credibility.
\end{itemize}

Overall, the structural composition features assess whether an article is well-built in terms of format and sourcing. These features reflect the Transparency dimension because they deal with how openly the article presents its information structure and sources to the reader. A well-structured article with clear sections, supporting images, and cited sources demonstrates transparency and professionalism. These are qualities we expect in higher quality categories. Additionally, by considering structural elements, we also indirectly support the Independence dimension: for example, using multiple sources and proper citations suggests independent verification of facts rather than just copying a single source. The combination of these structural features helps the model distinguish a thoroughly crafted article from a shoddy one on more than just vocabulary – it looks at the article’s journalistic form.

\subsection{Readability Metrics}

While the previous features focus on complexity and content, it is also important that a high-quality article is readable and accessible to its intended audience. To capture this, we include readability metrics that gauge how easy or difficult a text is to read. These features complement linguistic complexity by ensuring that we don’t just reward complexity for its own sake – the writing should also be clear. In other words, Attractiveness in journalistic quality is a balance between rich content and reader-friendly presentation. 

We adapt established readability formulas for the Czech language, accounting for Czech’s morphological complexity and free word order. Classic readability measures were originally developed for English, using factors such as sentence length and syllable count per word. However, applying them directly to Czech can be misleading because Czech words tend to be longer (due to inflectional endings) and the sentence structure can be more flexible. Therefore, we calibrate these formulas with Czech-specific considerations. Research has shown that a Russian-adapted formula works better for Czech than the raw English one, due to similarities in language structure, so we base our adaptation on similar principles (given the closeness of Slavic languages). In practice, this means our readability score might weight sentence length a bit more heavily (since very long sentences can be hard to parse in Czech) and adjust the baseline expectations for word length and syllable count. 

The outcome is a numeric readability score for each article. A higher score (or lower difficulty) indicates the article is easier to read, while a lower score means it’s more complex or dense. We expect high-quality articles to fall in a moderate range: they should not be trivial to read (as overly simplistic text might indicate shallow content), but they also should not be needlessly convoluted. By including the readability metric, we ensure the model can penalize articles that are linguistically complex but unreadable, and conversely recognize when a slightly simpler writing style still conveys high-quality information effectively. 

In summary, the readability feature provides a check on the Attractiveness dimension: it helps distinguish content that is both rich and well-written from content that is rich but poorly presented. This contributes to a more nuanced quality assessment, aligning our feature set with human perceptions of what makes an article good – not just what is written, but how it comes across to the reader.

\subsection{Metadata Features (Diversity)}
Articles don’t exist in isolation — contextual information about an article can also relate to its quality. In this context, Diversity refers to capturing a broad range of content conditions and ensuring the model is aware of different sources and patterns, so it can handle the variety present in the dataset.

Channel ID (Source Identifier): Each article in our dataset comes from a specific source or channel (for example, a particular news outlet or a section of the news platform). Initially, we included an anonymized channel identifier as a metadata feature. This allowed the model to recognize patterns in writing style or typical quality that might differ between various channels. Although this feature significantly improved classification performance, we ultimately decided not to use it in our final model.

The primary reason for excluding the channel identifier was to ensure the model classifies the articles based solely on textual content rather than relying on the reputation or historical patterns of specific channels. If we had used the channel ID, the model would strongly overfit to these identifiers—producing artificially high results by learning channel-specific biases. Consequently, if a particular channel aimed to improve its journalistic quality, the model might fail to reflect these improvements accurately, as it would still base predictions on outdated channel-specific patterns rather than on the actual article quality.

Furthermore, relying on channel identifiers could cause problems when classifying articles from channels that did not appear in the training dataset. In such cases, the model would lack critical information about these unseen channels, reducing its predictive reliability. By excluding this feature, we ensure the model remains generalizable and sensitive to the actual text content, facilitating a fair and unbiased assessment of article quality across all channels.
